\section{Discussion and limitations}
\label{sec:discussion}

\paragraph{What the hypothesis survived.} Coarse-to-fine strided windows get very close
to the optimum for the smooth, clustered layouts that motivate the problem. On the two
real 70{,}000-point UMAP embeddings, the default preset is within \RealBalancedExcess\ of
the certified optimum and visually indistinguishable from it. The method needs no dense
matrix, runs in 0.2--1.4\,s with all 24 threads at $10^5$ points (balanced preset), gives the same result for
any number of threads, and never increases the cost.

\paragraph{What it did not survive.} \ShuffleSnap\ is a local search, and its failures
follow from that. (i)~On near-uniform data, its relative excess grows with $N$ (for the
baseline, from about 1\% at $10^3$ points to 2.5\% at $10^5$ and 5.3\% at $10^6$). The excess consists of a
few long closed loops of one-cell shifts that no window can see whole. (ii)~Masks with
concave bays and grids with surplus cells need global rerouting of mass, which the
baseline gets wrong. With uniform points on the S-shaped mask it was 18\% above optimal,
worse than BSP-OT. The $\sqrt2$ schedule removes most of this (1.3\% for balanced, 0.13\%
for quality), but not all of it. (iii)~A good start does not help,
and structured starts hurt. (iv)~On a jittered lattice, where the optimum is almost the
identity, every heuristic we tested stays far from optimal in relative terms.

\paragraph{The exact solver changes the question.} We built the certified solver as an
evaluation tool. It turned out to be practical. Starting from \ShuffleSnap's answer, it proved optimality
at $10^5$ points in 7\,s to 23\,min on one core (median about 2\,min), and at $10^6$
uniform points in 12\,min.
When a certified optimum is worth the wait, it is the better choice at this scale, and
\ShuffleSnap's role is then that of a fast, good-enough answer, a warm start, and a
method for sizes and time budgets where exact solving is out of reach. We report this
even though it narrows the case for the method we set out to test.

\paragraph{Displacement is not neighbourhood.} BSP-OT reaches within a few percent of the
optimal displacement on the real embeddings, but halves the fraction of lattice
neighbours that were input neighbours, because merging many partitions shuffles points
locally. RasterFairy keeps local order inside its blocks and breaks it along block seams.
If neighbourhood preservation is the goal, it should be optimized directly (e.g.\ with
LAS/FLAS~\citep{barthel2023improved}); minimizing displacement buys it only to the
extent that the optimal transport map is smooth.

\paragraph{Limitations of this study.} Everything ran on one workstation, CPU only.
BSP-OT has a CUDA implementation that we did not test. The objective is squared
Euclidean displacement in 2D. The window solver works for any cost, but we did not
evaluate other costs or 3D grids. The certified solver depends on OR-Tools and on
integer scaling. Its certificate accounts for the scaling, but its running time depends
on the candidate set and the warm start. Real data consist of two UMAP embeddings with
one parameter setting. The bbox normalization is part of the problem definition; other
normalizations change the optimum. Timings include warm-up exclusions and a fixed
pinning policy, and unrelated background services were running (load about 1). We
compared against the representative settings of each comparator, but not every setting.
In particular, RasterFairy's swap optimizer was excluded because it requires dense
$N\times N$ memory and does not run on Python~3 as released.
